\begin{abstract}
Music recommenders learn what listeners like in general, but the song that fits
depends on the moment. Vibe Shuffle asks whether facial and cardiac changes
measured relative to personal references, together with recent camera movement,
can make next-song selection more context-sensitive. At the start of a session
the system initializes personal reference values for facial and cardiac signals. The facial reference then
adapts slowly to longer visible changes, whereas the cardiac reference remains
fixed. Later facial and cardiac observations are expressed as deviations from
the active references, while movement and nodding are summarized over a recent
camera window. The resulting observations are mapped to a valence--arousal plane
without assigning the listener an emotion label. Observations collected during a track form a
short trajectory whose mean identifies a catalog region, and the nearest
unplayed song in that region is selected. An exploratory within-session
crossover compared Vibe Shuffle with Random selection from the same catalog.
Mood-fit ratings were higher on average under Vibe Shuffle while general liking
differed little on average, but the confidence intervals for both differences
included zero, so no recommendation benefit is established. The contribution is
an inspectable method that turns within-session multimodal observations into a single
song-selection decision, together with a blueprint for the study needed to test
it.
\end{abstract}
